\subsection{Compute and memory of the router modules}
Tables~\ref{tab:mcp03-receipts}, \ref{tab:mcp02-pilot_gpu-receipts} and \ref{tab:mcp02-pilot_cpu-receipts} list, per module:
\begin{itemize}
\item measured parameters (total and without the hashed embedding);
\item recurrent depth in block applications;
\item FLOPs per scored candidate;
\item router latency per decision;
\item peak memory for training and inference.
\end{itemize}
All modules are small: the hashed embedding dominates the parameter count, and a scored decision costs
megaFLOPs to low gigaFLOPs. Recurrence trades parameters for FLOPs. At equal depth, a looped block or a TRM
holds the parameters of one block but pays the FLOPs of four, while the four-layer Transformer pays both.
Training memory is dominated by the framework runtime and the in-memory dataset rather than by the router. The
inference forward adds a few megabytes on top of the runtime.

\IfFileExists{tables/mcp03_receipts.tex}{\begin{table}[t]
\centering
\small
\begin{adjustbox}{max width=\linewidth}
\begin{tabular}{lrrrrrrrrrrrr}
\toprule
Router & Family & Params & Non-emb. & Depth & FLOPs/cand. & ms/dec. & Infer +MiB & Train CPU +MiB & Train CUDA MiB & s/step GPU & File R@2K+map & Sym R@2K \\
\midrule
dense\_d64 & dense & 301633 & 38977 & 0 & 66688 & 18.69 & 3 & 129 & 44 & 0.46 & 0.77 & 0.08 \\
hrm\_h2l1\_d64 & hrm & 368577 & 105921 & 4 & 34706688 & 162.09 & 23 & 756 & 645 & 0.49 & 0.78 & 0.12 \\
looped\_p4\_d64 & looped & 320513 & 57857 & 4 & 34706688 & 165.14 & 22 & 740 & 637 & 0.59 & 0.76 & 0.12 \\
looped\_p8\_d64 & looped & 320513 & 57857 & 8 & 69412096 & 308.11 & 22 & 1317 & 1204 & 0.96 & 0.76 & 0.17 \\
transformer\_l2\_d64 & transformer & 368449 & 105793 & 2 & 17353984 & 87.96 & 21 & 452 & 343 & 0.48 & 0.78 & 0.11 \\
transformer\_l2\_d64\_tokens & transformer & 368321 & 105665 & 2 & 17353984 & 85.20 & 20 & 452 & 343 & 0.47 & 0.77 & 0.10 \\
transformer\_l4\_d64 & transformer & 468417 & 205761 & 4 & 34706688 & 167.68 & 22 & 744 & 628 & 0.57 & 0.78 & 0.15 \\
trm\_c2i1\_d64 & trm & 318593 & 55937 & 4 & 34706688 & 167.06 & 22 & 740 & 644 & 0.52 & 0.75 & 0.13 \\
trm\_ldt\_c2i1\_d64 & trm\_ldt & 318918 & 56262 & 4 & 34707328 & 166.11 & 23 & 792 & 711 & 0.58 & 0.77 & 0.09 \\
\bottomrule
\end{tabular}
\end{adjustbox}
\caption{MCP-03 router modules: measured size, recurrent depth, FLOPs per candidate, latency, inference memory delta around a 32-candidate forward, training memory over 15 CPU / 40 GPU steps (CPU: growth above the loaded dataset and runtime; GPU: peak allocation), and held-out localisation in the LDT arm.}
\label{tab:mcp03-receipts}
\end{table}
}{}
\IfFileExists{tables/mcp02_pilot_gpu_receipts.tex}{\begin{table}[t]
\centering
\small
\begin{adjustbox}{max width=\linewidth}
\begin{tabular}{lrrrrrrrrrr}
\toprule
Router & Family & Params & Non-emb. & Depth & FLOPs/cand. & ms/decision & Infer RSS MiB & Train RSS MiB & Train CUDA MiB & Dev top-1 \\
\midrule
dense\_d128 & dense & 668801 & 143489 & 0 & 264448 & 15.54 & 358.10 & -- & -- & 0.84 \\
hrm\_h2l1\_d128 & hrm & 933761 & 408449 & 16 & 510429696 & 34.77 & 426.00 & 2470 & 927 & 0.92 \\
looped\_p12\_d128 & looped & 739329 & 214017 & 12 & 382822912 & 66.98 & 427.30 & 2446 & 2579 & 0.91 \\
looped\_p4\_d128 & looped & 739329 & 214017 & 4 & 127609344 & 36.22 & 427.90 & -- & 913 & 0.91 \\
sham\_shuffled\_transformer\_l2\_d128 & transformer & 933505 & 408193 & 2 & 63805952 & 20.84 & 427.30 & 2478 & 507 & 0.14 \\
transformer\_l1\_d128 & transformer & 735233 & 209921 & 1 & 31904256 & 15.91 & 427.10 & -- & -- & 0.91 \\
transformer\_l2\_d128 & transformer & 933505 & 408193 & 2 & 63805952 & 21.57 & 426.90 & -- & -- & 0.91 \\
transformer\_l2\_d128\_tokens & transformer & 933249 & 407937 & 2 & 63805952 & 21.69 & 427.60 & -- & -- & 0.73 \\
transformer\_l4\_d128 & transformer & 1330049 & 804737 & 4 & 127609344 & 30.53 & 427.00 & -- & -- & 0.91 \\
trm\_c2i1\_d128 & trm & 735489 & 210177 & 4 & 127609344 & 38.14 & 427.50 & -- & 924 & 0.91 \\
trm\_c4i2\_d128 & trm & 735489 & 210177 & 12 & 382822912 & 67.73 & 429.30 & 2458 & 2595 & 0.91 \\
trm\_ldt\_c2i1\_d128 & trm\_ldt & 736134 & 210822 & 4 & 127610624 & 37.38 & 427.40 & 2498 & 922 & 0.90 \\
\bottomrule
\end{tabular}
\end{adjustbox}
\caption{MCP-02 pilot gpu router receipts: measured parameters, recurrent depth (block applications), FLOPs per candidate (FlopCounterMode), router latency, peak memory.}
\label{tab:mcp02-pilot_gpu-receipts}
\end{table}
}{}
\IfFileExists{tables/mcp02_pilot_cpu_receipts.tex}{\begin{table}[t]
\centering
\small
\begin{adjustbox}{max width=\linewidth}
\begin{tabular}{lrrrrrrrrrr}
\toprule
Router & Family & Params & Non-emb. & Depth & FLOPs/cand. & ms/decision & Infer RSS MiB & Train RSS MiB & Train CUDA MiB & Dev top-1 \\
\midrule
dense\_d48 & dense & 219985 & 22993 & 0 & 37728 & 22.54 & 3.30 & -- & -- & 0.48 \\
hrm\_h2l1\_d48 & hrm & 257905 & 60913 & 16 & 82293696 & 96.33 & 22.00 & 1482 & -- & 0.61 \\
looped\_p4\_d48 & looped & 231073 & 34081 & 4 & 20574144 & 96.91 & 20.60 & -- & -- & 0.56 \\
sham\_shuffled\_transformer\_l2\_d48 & transformer & 257809 & 60817 & 2 & 10287552 & 49.17 & 20.90 & 1313 & -- & 0.16 \\
transformer\_l1\_d48 & transformer & 229537 & 32545 & 1 & 5144256 & 49.61 & 19.00 & -- & -- & 0.56 \\
transformer\_l2\_d48 & transformer & 257809 & 60817 & 2 & 10287552 & 81.31 & 20.60 & -- & -- & 0.62 \\
transformer\_l2\_d48\_match & transformer & 257905 & 60913 & 2 & 10287552 & 49.30 & 20.30 & 1290 & -- & 0.88 \\
transformer\_l4\_d48 & transformer & 314353 & 117361 & 4 & 20574144 & 84.70 & 20.90 & -- & -- & 0.65 \\
trm\_c2i1\_d48 & trm & 229633 & 32641 & 4 & 20574144 & 366.54 & 21.30 & -- & -- & 0.61 \\
trm\_c2i1\_d48\_match & trm & 229729 & 32737 & 4 & 20574144 & 107.48 & 22.10 & 1512 & -- & 0.86 \\
trm\_ldt\_c2i1\_d48 & trm\_ldt & 229878 & 32886 & 4 & 20574624 & 100.19 & 22.10 & 1483 & -- & 0.61 \\
\bottomrule
\end{tabular}
\end{adjustbox}
\caption{MCP-02 pilot cpu router receipts: measured parameters, recurrent depth (block applications), FLOPs per candidate (FlopCounterMode), router latency, peak memory.}
\label{tab:mcp02-pilot_cpu-receipts}
\end{table}
}{}
